\setcounter{section}{2}
\section{Algorithm for \Implication}
\label{sec:implalgo}
%
In this section, we give an algorithm for the problem $\Implication = \SAT(\{\Impl\})$ that is much faster than brute force and achieves $O(\sqrt{k})$ dependence of $k$ in the exponent $n$. 
For convenience, we reduce \Implication to the following problem. (Recall that for any graph $G=(V,E)$, we say that $S\subseteq V$ is \emph{closed}, if for all $(u,v)\in E$, we have $u \in S$ implies $v\in S$.)
%
\begin{problem}[\wDAGImplication]
Given an DAG $G=(V,E)$ with node weights $w: V\to \nats$ and parameter $k\in \nats$, determine whether there is a closed set $S \subseteq V$ of weight exactly $k$, i.e., $w(S) = k$.
\end{problem}
The easy reduction works as follows. For each variable $x_i$, we introduce a corresponding vertex $x_i$ of weight 1 and introduce an edge $(x_i, x_j)$ for every implication constraint $x_i \Rightarrow x_j$ of $\phi$. We contract each strongly connected component $C=\{v_1, v_2, \dots, v_\ell\}$ in $G$ to a single vertex $v_C$ of weight $\sum_{i=1}^\ell w(v_i)$ in time $\Oh(n+m)= \Oh(n^2)$~\cite{Tarjan72}. Observe that the resulting graph is a DAG which has a closed set of weight~$k$ if and only if $\phi$ has satisfying assignment of weight $k$.
%
Recall that for any $v\in V$, we let $D(v)$ denote the set of descendants of $v$, i.e., the set of nodes reachable from $v$ (including $v$).
%
As we will formally argue later, by a $f(k)n^{\Oh(\sqrt{k})}$-time preprocessing it is not difficult to preprocess a \wDAGImplication instance into the following form, which we call \emph{Frobenius instance}, as it admits a combinatorial characterization of solvability that is analogous to Schur's bound for the Frobenius coin problem.
%
\begin{defn}
A \emph{Frobenius instance with parameter $k$} is a weighted directed graph $G = (V,E, w)$ with $\ell$ parts $V = V_1 \cup V_2 \cup \cdots \cup V_\ell$ and weight function $w:V \to \nats$ such that the following properties hold: 
\begin{enumerate}[label=(P\arabic*)]
\item\label{item:structweights} there are weights $w_1, \dots, w_\ell$ such that $w(v) = w_i$ for all $v\in V_i$ and $i \in [\ell]$.
\item\label{item:structedges} for any edge $(u,v)\in E$, we have $u\in V_i$ and $v\in V_j$ for some $\ell \ge i > j \ge 1$,
\item\label{item:size} for all $i\in [\ell]$, we have $|V_i| \ge k$, 
\item\label{item:weights} for all $v\in V$, we have $w(D(v)) \le \sqrt{k/2}$.
\end{enumerate}
\end{defn}
%
Intuitively, the necessary preprocessing follows from the following arguments: To ensure~\ref{item:weights}, note that any weight-$k$ closed set $S$ has at most $\sqrt{2k}$ many vertices $v\in S$ with $w(D(v)) > \sqrt{k/2}$, which we can exhaustively enumerate with $n^{\Oh(\sqrt{k})}$-time overhead. By suitably arranging remaining nodes among the layers, it is straightforward to ensure~\ref{item:structweights}, \ref{item:structedges} and additionally that $\ell \le f(k)$, since by~\ref{item:weights}, each node has at most $\Oh(\sqrt{k})$ descendants. Finally, to ensure~\ref{item:size}, if any part $V_i$ is small (i.e., $|V_i| < k$), we can exhaustively try out including any subset of $V_i$, introducing an overhead of only $2^{\Oh(k)}$ per $V_i$; since $\ell \le f(k)$, this additional overhead is bounded by $f(k)2^{\Oh(k)}$.  
%
%
%
%
%
If a Frobenius instance had no edges, then Schur's bound on the Frobenius coin problem implies that it has a solution if and only if $\gcd(w_1,\dots, w_\ell) \divides k$. We prove that this criterion holds even in the setting of precedence constraints.  
%
\begin{lem}\label{lem:Frobcrit}
Let $G$ be a Frobenius instance with  parameter $k$. Then $G$ has a closed set of weight~$k$ if and only if $\gcd(w_1,\dots, w_\ell) \divides k$.
\end{lem}
\begin{proof}
Since $\gcd(w_1,\dots, w_\ell) \divides w(S)$ for any $S\subseteq V$, the condition $\gcd(w_1,\dots,w_\ell) \divides k$ is necessary for $G$ to have a closed set of weight $k$.
%
We show that this condition is also sufficient via induction on $\ell$. In the base case $\ell=1$, let $S \subseteq V_1=V$ be an arbitrary subset of $k/w_1$ vertices (note that by $k/w_1 \le k \le |V_1|$, such a set indeed exists). By construction, $S$ has weight $|S|w_1 = k$ and is closed, as $G$ cannot contain any edges.
%
Thus let us assume that the claim holds for all $\ell' \le \ell - 1$ and consider a Frobenius instance with $d':=\gcd(w_1,\dots, w_\ell) \divides k$. Let $d := \gcd(w_1,\dots,w_{\ell-1})$. We may assume that $d \notdivides k$; otherwise, already the Frobenius instance $G[V_1 \cup \dots \cup V_{\ell-1}]$ satisfies the assumption $\gcd(w_1,\dots, w_{\ell-1})\divides k$ and we obtain a closed set by inductive hypothesis.
%
Intuitively, we want to use the variables in $V_\ell$ to reach the target weight $k$ modulo $d$; then we can reduce to a simpler instance where every weight (including the target weight) is divided by $d$.
Note that we may assume 
\begin{equation}\label{eq:boundsond}
 2 \le d \le \sqrt{k/2},
 \end{equation}
where the lower bound follows from $d\notdivides k$ and the upper bound follows from $d\le \min_{i\in [\ell-1]} w_i \le \sqrt{k/2}$, as $w(v) \le w(D(v)) \le \sqrt{k/2}$ for any $v\in V$.
%
Let $b$ be the smallest non-negative integer such that $b \cdot  w_\ell \equiv k \pmod d$. Such an integer exists and satisfies $b < d$: By B\'ezout's identity, since $\gcd(w_\ell, d) = d' \divides k$, there are coefficients $\beta, \gamma$ such that $\beta w_\ell + \gamma d = k$, and thus any $b$ with $b \equiv \beta \pmod d$ achieves the desired congruence. 
%
Let $S\subseteq V_\ell$ be an arbitrary subset of size $b< d$; such a set indeed exists as $d\le \sqrt{k/2} \le k \le |V_\ell|$. We observe that $S$ satisfies
\begin{equation}\label{eq:boundonwDs}
w(D(S)) \le \sum_{s \in S} w(D(s)) \le |S| \sqrt{k/2} \le d\sqrt{k/2} \le \frac{k}{2},
\end{equation}
where we used~\ref{item:weights} for the second inequality, and~\eqref{eq:boundsond} for the last inequality. 
Consider the graph $G' = (V', E')$ obtained as a copy from $G$ from which we delete $V_\ell \cup D(S)$ and define the node weights $w'(v') = w(v)/d$ for any $v\in V\setminus (V_\ell \cup D(S))$. We claim that $G'$ is a Frobenius instance with parameter $k' := (k-w(D(S)))/d$ (observe that $k'$ is indeed integer, as $w(D(S)) \equiv b w_\ell \equiv k \pmod d$, and that $k'\ge 0$ by~\eqref{eq:boundonwDs}). If this is indeed the case, then by inductive hypothesis $G'$ has a closed set $S'$ with $w'(S') = k'$, since the gcd of the weights $w'$ is 1. Observe that by construction, $D(S)\cup S'$ is a closed set in $G$ of weight $w(D(S)) + d \cdot w'(S') = w(D(S)) + (k-w(D(S))) = k$, as desired.
%
It remains to prove that $G'$ is indeed a Frobenius instance with parameter $k'$. First, observe $G'$ has $\ell-1$ layers $V'_i := V_i \setminus D(S), i\in [\ell-1]$ and that $w'$ is well defined, as $d\divides w_i$ for all $i\in [\ell-1]$. Conditions~\ref{item:structweights} and~\ref{item:structedges} of being Frobenius are fulfilled as $G'$ is a subgraph of $G$. To see~\ref{item:size}, note that 
\[ |V'_i| \ge |V_i| - |D(S)| \ge |V_i|- w(D(S)) \ge k-w(D(S)) \ge k'. \]
To see~\ref{item:weights}, we observe that by~\eqref{eq:boundonwDs} and~\eqref{eq:boundsond}, we have
\[ k' = \frac{k-w(D(S))}{d} \ge \frac{k-k/2}{d} = \frac{k}{2d} \ge \frac{k}{d^2}. \]
Thus, for any $v'\in V'$, we obtain
\[ w'(D(v')) \le \frac{w(D(v))}{d} \le \frac{\sqrt{k/2}}{d} = \sqrt{\frac{k}{2d^2}} \le \sqrt{k'/2}, \]
where we used condition~\ref{item:weights} of $G$ in the second inequality. Thus, $G'$ is indeed a Frobenius instance with parameter $k'$, concluding the claim and thus the proof of our lemma.
\end{proof}
%
%
The above criterion is the main technical tool in the algorithmic result of the session. What remains is to show that the instance can be preprocessed in a way that it becomes a Frobenius instance.
\begin{thm}\label{thm:implalgo}
We can solve \wDAGImplication in time $f(k)n^{4\sqrt{k}}$.
\end{thm}
\begin{proof}
Consider the following recursive algorithm, which proceeds in 4 steps:
%
\emph{Step 1:} For every $v\in V$ with $w(D(v)) \ge \sqrt{k/2}$, we return YES if a recursive call determines that $G[V \setminus D(v)]$ has a closed set of weight $k-w(D(v))$; otherwise, we delete $v$ and all its ancestors from $G$. From now on, $G$ satisfies $w(D(v)) \le \sqrt{k/2}$ for all $v\in V$.
%
\emph{Step 2:} We construct layers $L_1, \dots, L_{\sqrt{k/2}}$ by the following iterative process: for every $i=1, \dots, \sqrt{k/2}$, we let $L_i$ consists of all vertices in $V\setminus (L_1 \cup \cdots \cup L_{i-1})$ whose outgoing edges end in $L_1\cup \cdots \cup L_{i-1}$. Note that $L_1, \dots, L_{\sqrt{k/2}}$ partitions $V$; in particular, every vertex is included in some $L_i$, since if there was a vertex $v\in V\setminus (L_1 \cup \cdots \cup L_{\sqrt{k/2}})$, then by construction there exists a path from $v$ containing strictly more than $\sqrt{k/2}$ vertices, leading to the contradiction $w(D(v))\ge |D(v)| > \sqrt{k/2}$. 
%
We observe that each layer $L_i$ can be partitioned into sublayers $L_{i,j}, j\in \{1, \dots, \sqrt{k/2}\}$ such that each $v\in L_{i,j}$ has weight $w(v) = j$: there can be no vertex of larger weight, as otherwise $w(D(v)) \ge w(v) > \sqrt{k/2}$ yields a contradiction. We consider layers $L_{i,j}$ in increasing lexicographic order of $(i,j)$: If $|L_{i,j}| < k$, then for every $v\in L_{i,j}$, we return YES if a recursive call determines that $G[V\setminus D(v)]$ contains a closed set of size $k-w(D(v))$, and otherwise we delete $v$ and all its ancestors from~$G$. Observe that by the lexicographic ordering, we never delete vertices from already processed layers, so that at the end of the process, each $L_{i,j}$ is either empty or contains at least $k$ vertices. 
%
\emph{Step 3:} We let $V_1, \dots, V_\ell$ be an enumeration of all non-empty sublayers $L_{i,j}$ by the lexicographic order on $(i,j)$ so that any vertex $v\in V_{i}$ has only edges to vertices in $V_1 \cup \cdots \cup V_{i-1}$. Observe that by construction, this yields a Frobenius instance. Let $w_1,\dots, w_\ell$ be the weights of the Frobenius instance. We return YES if $\gcd(w_1,\dots, w_\ell) \divides k$ and NO otherwise. 
%
Using Lemma~\ref{lem:Frobcrit}, the correctness of the algorithm is easy to see.
%
\begin{claim}\label{cl:imp1}
The above algorithm is correct.
\end{claim}
\begin{proof}
If the algorithm returns YES, indeed there is a closed set of size $k$: 
If we return YES in Steps 1 or 2, we have found a vertex $v$ and a closed set $S'$ in $G[V\setminus D(v)]$ of size $k-w(D(v))$, which yields a closed set $S'\cup D(v)$ in $G$ of size $k$, as desired. Otherwise, we have arrived at a Frobenius instance and returned YES since $\gcd(w_1,\dots, w_\ell)\divides k$, which implies that $G$ has a closed set of size $k$ by Lemma~\ref{lem:Frobcrit}.
%
Conversely, fix a closed set $S$ of size $k$, and we show that the algorithm returns YES:
If $S$ contains a vertex $v$ investigated in Steps 1 or 2, then the recursive call to $G[V\setminus D(v)]$ (for the first such vertex $v$) will find a solution of size $|S|-w(D(v))$ (note that $D(v) \subseteq S$ if $v\in S$). Otherwise, we have arrived at a Frobenius instance which must satisfy $\gcd(w_1,\dots, w_\ell)$ by Lemma~\ref{lem:Frobcrit}, and we return YES.
\cqed\end{proof}
%
Finally, we need to bound the running time of the recursive algorithm. The analysis relies on the observation that the algorithm makes at most $n$ recursive calls with a parameter decrease of at least $\sqrt{k/2}$, and at most $O(k^2)$ recursive calls with a parameter decrease of one.
\begin{claim}\label{cl:imp2}
The above algorithm can be implemented in time $f(k) n^{4\sqrt{k}}$.
\end{claim}
\begin{proof}
Let $U$ be the set of vertices of small layers ($|L_{i,j}| < k$) considered in Step 2. We observe that the above algorithm can be implemented recursively with the following recurrence on its running time $T(n,k)$ on instances with $n$ vertices and parameter $k$.
\begin{align*}
T(n,k) & \le \sum_{v\in V, w(D(v)) \ge \sqrt{k/2}} T(n, k-w(D(v))) + \sum_{u \in U} T(n,k-w(D(u))) + \Oh(n^2) 
\end{align*}
We claim by induction on $k$ that this yields a bound of $T(n,k) \le f(k) n^{4\sqrt{k}}$ for some $f(k) = k^{\Oh(k)}$. It is not difficult to see that for $k\le 2$, we can solve the problem in time $\Oh(n^2) = \Oh(n^{4\sqrt{k}})$, yielding the base case. For $k\ge 3$, we thus obtain the following bound, using that in Step 2, we process less than $k$ vertices for each ``small'' sublayer $L_{i,j}, 1\le i,j \le \sqrt{k/2}$, i.e., $|U|\le k (k/2) = k^2/2$,
\begin{align*}
T(n,k) & \le \Oh(n \cdot f(k-\sqrt{k/2}) n^{4\sqrt{k-\sqrt{k/2}}} + k^2 f(k-1) n^{4\sqrt{k-1}} + n^2) \\
   & \le (f(k)/2) (n^{4\sqrt{k-\sqrt{k/2}}+1} + n^{4\sqrt{k}}) \le f(k)n^{4\sqrt{k}},
  \end{align*}
 where the second bound follows from choosing $f(k) = k^{\Oh(k)}$ large enough to ensure $k^2f(k-1) \le f(k)/2$ and 
 the last bound follows from the observation that $4\sqrt{k-\sqrt{k/2}} + 1\le 4\sqrt{k}$ if and only if 
   \begin{align*}
& &  \left(4\sqrt{k-\sqrt{k/2}} + 1\right)^2 & \le 16k\\
& \iff & 16(k-\sqrt{k/2}) + 8\sqrt{k-\sqrt{k/2}} + 1 & \le 16k\\
& \iff & 8\sqrt{k-\sqrt{k/2}} + 1 & \le 16\sqrt{k/2},
  \end{align*}
  where the last inequality holds since $8\sqrt{k}+1 \le 16\sqrt{k/2}$ as $k\ge 3$.
  \cqed\end{proof}
Claims \ref{cl:imp1} and \ref{cl:imp2} show the correctness of our algorithm for \wDAGImplication. By the reduction described at the beginning of the section, a similar algorithm follows for \Implication.
\end{proof}
